\chapter{Methodology}
\label{chap:methodology}

%% PURPOSE: This chapter explains the HOW of the research — the pipeline from raw
%% data to validated choreographies. It must justify every design choice and connect
%% each stage explicitly to an RQ. The synthetic-data section is load-bearing: it
%% must convince the reader that the pipeline can be trusted before the real
%% Pearson data is analysed.

\section{Overview}
\label{sec:methodology-overview}
% TODO: Write a short (3-5 sentence) roadmap of the chapter. Describe the four
% pipeline stages (synthetic data generation → feature engineering → pattern
% discovery → evaluation) and state which RQ each stage is primarily designed
% to address.
This chapter describes the methodological process used to analyze, characterize, and predict the academic success of students in fully online schooling.
The pipeline comprises these stages: data extraction and preparation, pattern discovery through k-means clustering and UMAP visualization, association of the discovered patterns with academic outcomes, and prediction of student success using a Random Forest model.

\section{Synthetic Data as a Methodological Choice}
\label{sec:synthetic-data-rationale}
% TODO: Justify why synthetic data was used as the primary substrate for this
% dissertation, covering: ...
The data used in this dissertation is synthetic. The real datasets were expected first from Pearson and subsequently from Stride Inc., but neither could be delivered in time. The synthetic dataset had originally been created as a working substrate to allow development to proceed while the real data was pending; once it became clear the real data would not arrive, continuing with the synthetic data was the only viable path.

This choice carries a significant methodological advantage. Because the dataset was generated by the author, the true behavioral archetype of each student is known by construction. This ground truth, which is never exposed to the clustering algorithm, makes it possible to measure whether the discovered clusters actually recover the underlying behaviors, using metrics such as the Adjusted Rand Index (ARI). With real, anonymized data, these true labels would be unknown, and such a quantitative validation of the clustering would not be possible.

A second advantage is openness. As the data does not belong to real students or institutions and was created entirely by the author, it can be used and published without privacy or licensing constraints, supporting full reproducibility of the pipeline.

This approach has an explicit limitation. Synthetic results cannot be validated in the same way as real data, and they do not constitute empirical claims about real Connections Academy students. What they validate is the methodology, the pipeline and analytical methods developed to carry out the study.

The generator's parameters were designed through informed intuition rather than calibrated against measured empirical distributions. The five archetypes were specified so that each would be behaviorally plausible and distinct from the others, drawing on the patterns of engagement and disengagement discussed in the literature reviewed earlier~\cite{fostering_curtis_2015, beck_atrisk_2024}. They are not intended to reproduce exact empirical frequencies from real online schools, but to provide coherent, recognizable behavioral profiles against which the discovery pipeline can be developed and validated.

\section{Synthetic Data Generation}
\label{sec:synthetic-data-generation}
% TODO: Describe the synthetic data generator in detail: ...
The synthetic generator builds an academic year of activity for thousands of students, each assigned to one of five behavioral archetypes that describe how they study. The Steady Worker is a consistent, engaged student, focused on content study, assignments, and assessments. The Balanced Engager is the most social and interactive profile, driven by live sessions and discussion forums. The Night Crammer studies intensively, but concentrates activity late at night or in the early morning. The Minimal Browser mostly navigates pages, checks the dashboard, and looks up grades, with little actual studying. Finally, the Disengaged Ghost is a near-absent student, showing minimal activity of any kind.

Students are sampled from these archetypes with embedded correlations. At-risk students are more likely to be assigned to the lower-engagement archetypes, reflecting the weaker outcomes associated with these profiles. The dataset reproduces a realistic institutional structure: 32 schools across 32 U.S. states, spanning 13 grade levels from Kindergarten to Grade 12 (K--12).

Activity is generated as event-level logs, one student at a time. Each event records an action type, drawn according to the student's archetype, along with the time of day and, where applicable, the session duration. A fixed random seed (42) is used throughout the generation process to ensure full reproducibility.

\section{Feature Engineering}
\label{sec:feature-engineering}
% TODO: Describe the weekly feature matrix construction: ...
While the theoretical concept of a learning choreography implies a sequence of actions over time, this methodology practically operationalizes it through the aggregation of weekly behavior patterns and proportions. The unit of analysis is the student-week: one row per student per calendar week. This choice provides enough data points for the clustering algorithm to find patterns, while letting each student be represented as a sequence of weeks across the year. Two alternatives were considered, the student-day and the student-year, but the former would generate an excessive volume of sparse data, and the latter would not retain enough observations per student. The student-week was the most balanced option.

Weekly features follow a two-block design. The WHAT block captures the proportion of each semantic action type performed during the week (for example, how much content study, assessment, or forum activity). The WHEN block independently captures the proportion of activity falling in each of four time-of-day bins: Morning, Afternoon, Evening, and Night.

The two blocks are kept separate rather than crossed. In an early attempt, action and time-of-day were combined into a single set of features, which caused the time dimension to dominate k-means: the resulting clusters separated students almost entirely by period of day rather than by behavior. Decoupling the two blocks resolved this and let the clustering focus on what students do, with time of day as a secondary signal.

Login and Logout actions are deliberately excluded from the WHAT block. They are universal, every student performs them, and therefore carry no discriminative value; they merely mark session boundaries. Finally, aggregate metadata (number of sessions, active days, total actions) is computed but kept outside the clustering feature space, so that students are grouped by the shape of their behavior rather than by sheer activity volume.

\sectionsection{Pattern Discovery (Clustering)}
\label{sec:pattern-discovery}
% TODO: Describe the clustering pipeline ...
At this stage, the goal of clustering is to discover the shape of behavior, for which k-means was chosen as a clustering algorithm widely used in educational data mining~\cite{romero_edm_survey}. Because k-means groups data by computing mathematical distances between points, raw action counts would let differences in scale dominate: the algorithm would simply separate students who click a lot from students who click little, failing to capture the actual behavioral patterns. To prevent this, each weekly vector is L1-normalized into proportions. After normalization, two students with the same behavioral mix, for instance, both spending half their activity on content study and half on forums, are treated as similar even if one performed many more actions than the other. Clustering therefore focuses on the composition of behavior rather than its volume.

Standardization is then applied so that each feature contributes comparably to the distance computation, preventing frequent actions from overwhelming rarer but discriminative ones. This step, however, required care. Some action types are very rare, almost no student performs them, so their columns have a near-zero standard deviation. Standardizing divides by the standard deviation, so these columns would trigger division-by-zero errors that corrupt the model. The \texttt{safe\_standardize} function avoids this by detecting near-constant columns (standard deviation below a small threshold) and neutralizing them to zero rather than dividing by them, so they simply carry no weight in the clustering.

The number of choreographies, $k$, was selected by evaluating solutions across a range of $k=2$ to $10$, computing both the inertia (Elbow method) and the Silhouette score. These automatic metrics did not, on their own, give a definitive answer. The Silhouette score favored sharp mathematical partitions, but the UMAP projection showed that behaviors merged continuously rather than forming isolated islands; relying on the metric alone would have split genuine profiles in half. Both $k=4$ and $k=5$ were therefore examined directly. Validation against the ground-truth archetypes showed that the two solutions recovered the original behaviors with essentially the same accuracy (ARI $\approx 0.53$). The deciding factor was interpretability: at $k=5$, two of the clusters were both low-engagement profiles, separated only by time of day, which made them redundant. The final choice of $k=4$ retains the same fidelity to the underlying structure while producing distinct, non-redundant choreographies.

The final clustering was produced by fitting k-means with $k=4$, ten initializations (\texttt{n\_init} $= 10$), and a fixed random seed (42) to guarantee reproducible results. For visual diagnosis, the resulting clusters were projected into two dimensions using Uniform Manifold Approximation and Projection (UMAP), with \texttt{n\_neighbors} $= 15$ and \texttt{min\_dist} $= 0.1$. It is important to stress that UMAP was used strictly for human inspection: it allowed visual verification of whether the discovered groups separated distinctly, but had no influence whatsoever on the cluster assignments, which were determined solely by k-means. Removing the projection would leave the clusters unchanged.

\section{Evaluation Design}
\label{sec:evaluation-design}
% TODO: Map each research question to its evaluation method and metric(s)...
RQ1 (Identification) was evaluated by comparing the algorithm's cluster assignments against the generator's ground truth. A contingency table (cluster $\times$ archetype) was built, and the purity of each cluster was computed, identifying the dominant student profile within each discovered group. To assess this correspondence quantitatively, two agreement metrics were applied: the Adjusted Rand Index (ARI) and Normalized Mutual Information (NMI). These metrics quantify the extent to which the algorithm recovered the behavioral structure originally embedded in the synthetic dataset.

RQ2 (Robustness) was assessed along two dimensions. Temporal stability was measured over the full academic year using week-over-week ARI and the percentage of students retained in the same cluster from one week to the next. A transition matrix was built to map how students moved between choreographies; the concentration of values along the main diagonal indicates the degree to which students maintained consistent learning routines rather than shifting abruptly between profiles. Cross-context stability examined whether the discovered patterns were universal rather than demographic artefacts. The percentage distribution of each choreography was evaluated across grade levels and states, checking that the standard deviation across contexts remained small. Finally, to assess whether the shape of behavior changed with schooling level, the cosine similarity of each profile's behavioral signature was computed across grade bands (K--5 versus 6--12), measuring how closely the underlying structure of these profiles is preserved.

RQ3 (Association with Outcomes) was assessed by first identifying each student's dominant choreography, the profile in which they spent most of their weeks across the year. With students grouped by dominant profile, a one-way ANOVA was applied to their final grades to test whether the performance differences between profiles were statistically significant, and the eta-squared statistic was used to measure the magnitude of the effect. Finally, to make the comparison fairer, these results were stratified by the student's initial at-risk status. This made it possible to examine to what extent final grades reflected the learning routines themselves rather than difficulties the student already carried beforehand.

RQ4 (Predictive Utility) was evaluated by training a Random Forest classifier to predict whether a student would pass, using features drawn only from the first 4 and 8 weeks of the year. This strict temporal cutoff prevents data leakage: it ensures the model cannot "peek" at end-of-year behavior to anticipate the outcome. Two models were compared using the ROC-AUC metric: a baseline using only activity-volume features, and an experimental model adding the student's dominant choreography. ROC-AUC was chosen over accuracy because the classes are imbalanced (most students pass), which would make accuracy misleading.

RQ5 (Actionability) was addressed by analyzing the trained prediction model to identify which actions carried the most weight in the final outcome, measuring feature importance in two ways: Gini importance and permutation importance. Two approaches were used because Gini importance has a well-known bias: it tends to overvalue features with many distinct values, such as navigation clicks, which can be misleading. Permutation importance is more reliable, as it measures the actual drop in model performance when a feature is shuffled. Once the most important behaviors were identified, the direction of effect was examined, checking whether a high proportion of a given action was associated with passing or failing, which yields the practical, actionable insights about which behaviors to encourage. Mean absolute SHAP values (via a TreeExplainer) were also considered as a means of corroborating these rankings.

\section{Reproducibility and Implementation Details}
\label{sec:reproducibility}
The pipeline can be reproduced in full, from start to finish. The guiding rule was to fix the random seed (42) at every stage of the project, from the generation of the synthetic students to the train/test split in the Random Forest. This forces the computation to make exactly the same decisions on any machine, guaranteeing identical results across runs. The analysis is organized as a sequence of Python scripts that must be executed in a defined order: synthetic data generation, feature engineering, clustering, validation, stability analysis, outcome association, early prediction, and interpretation, documented in the repository. The entire work was made available in a public repository, relying only on standard Python libraries (pandas, numpy, scikit-learn, umap-learn, matplotlib, seaborn, scipy).